\noindent\unskip\SourceLeaf{19}with the expression congruent to either 1, $\rho$, or $\rho^2$ modulo $q$, where $\rho$ denotes the root of unity
$\dfrac{-1+i\sqrt{3}}{2}$. Now, since in the domain of complex integers formed from $\rho$ the solutions
of the congruence $x^{q-1}\equiv 1\;(\mathrm{mod.}\ q)$ are the $(q-1)$ rational integer residues distinct from 0 modulo $q$,
the division into three cases made by \textsc{Gauss} by means of the congruences $B\equiv 0$, $B\equiv A$, $B\equiv -A\;(\mathrm{mod.}\ q)$ amounts to saying that the cubic character $\left[\dfrac{\pi_1}{q}\right]$ equals 1, $\rho$, or $\rho^2$, respectively. In the other case $q=3n+1$,
the simplest expression of the reciprocity law has yet to be reached here, because the factorization of
$q$ into its prime divisors has yet to be carried out. The connection between \textsc{Gauss}'s theorem and the reciprocity law is as follows.

If the factorizations of the rational primes $p, q$ into their primary prime factors are given by
\[ p=\pi_1.\pi_2,\qquad q=\varkappa_1.\varkappa_2,\qquad \pi_k'=2\pi_k=a\pm 3b\sqrt{-3}, \]
then one must set
\[
\left.\begin{array}{l}
\left[\dfrac{\pi_1'}{\varkappa_1}\right]=\rho^{\varepsilon_1}\equiv {\pi_1'}^{\frac{q-1}{3}}\equiv A+B\rho-B\rho^2 \\[2mm]
\left[\dfrac{\pi_2'}{\varkappa_1}\right]=\rho^{\varepsilon_2}\equiv {\pi_2'}^{\frac{q-1}{3}}\equiv A-B\rho+B\rho^2
\end{array}\right\}(\mathrm{mod.}\ \varkappa_1).
\]

\noindent \textsc{Gauss}'s division into three cases then agrees with the following division
\[
\text{I.}\quad \varepsilon_1\equiv \varepsilon_2 (\mathrm{mod.}\ 3),\qquad \left[\dfrac{\pi_1'}{\varkappa_1}\right]=\left[\dfrac{\pi_2'}{\varkappa_1}\right],
\]
\[
\text{II.}\quad \varepsilon_1\equiv \varepsilon_2+1\;(\mathrm{mod.}\ 3),\qquad \left[\dfrac{\pi_1'}{\varkappa_1}\right]=\rho\left[\dfrac{\pi_2'}{\varkappa_1}\right],
\]
\[
\text{III.}\quad \varepsilon_1\equiv \varepsilon_2-1\;(\mathrm{mod.}\ 3),\qquad \left[\dfrac{\pi_1'}{\varkappa_1}\right]=\rho^2\left[\dfrac{\pi_2'}{\varkappa_1}\right],
\]
where in the last three equations $\pi_k$ can evidently also take the place of $\pi_k'$. Indeed, if e.\,g.
case II holds, then adding and subtracting the above congruences with respect to the modulus $\varkappa_1$ gives
\[
2A\equiv \rho^{\varepsilon_2}(1+\rho)\equiv -\rho^{\varepsilon_2-1},\qquad 2B\equiv -\rho^{\varepsilon_2-1}\qquad\;(\mathrm{mod.}\ \varkappa_1),
\]
so that $(A-B)$ is divisible by $\varkappa_1$ and hence, being a rational number, also by $q$. Thus we do indeed have $B\equiv A\;(\mathrm{mod.}\ q)$. The reciprocity law now gives
\[
\left[\dfrac{\pi_1'}{\varkappa_1}\right]=\left[\dfrac{2}{\varkappa_1}\right].\left[\dfrac{\pi_1}{\varkappa_1}\right]=\left[\dfrac{2}{\varkappa_1}\right].\left[\dfrac{\varkappa_1}{\pi_1}\right],
\]
\[
\left[\dfrac{\pi_2'}{\varkappa_1}\right]=\left[\dfrac{2}{\varkappa_1}\right].\left[\dfrac{\pi_2}{\varkappa_1}\right]=\left[\dfrac{2}{\varkappa_1}\right].\left[\dfrac{\varkappa_1}{\pi_2}\right],
\]
so that in case II we further obtain
\[
\left[\dfrac{\varkappa_1}{\pi_1}\right]=\rho\left[\dfrac{\varkappa_1}{\pi_2}\right]=\rho\left[\dfrac{\varkappa_2}{\pi_1}\right]^2,\qquad
\left[\dfrac{\varkappa_1}{\pi_1}\right].\left[\dfrac{\varkappa_2}{\pi_1}\right]=\left[\dfrac{q}{\pi_1}\right]=\rho,
\]
\[
\left[\dfrac{q}{\pi_1}\right]\equiv q^{\frac{p-1}{3}}\equiv -\tfrac12+\tfrac12\sqrt{-3}\equiv -\tfrac12-\tfrac12\dfrac{a}{3b}\qquad\;(\mathrm{mod.}\ \pi_1).
\]
But if the rational number satisfies $q^{\frac{p-1}{3}}+\tfrac12+\tfrac12\dfrac{a}{3b}\equiv 0 \mathrm{mod.}\ \pi_1$, then it is also divisible by $\pi_2$ and hence by $p$.
This gives the congruence stated under II by \textsc{Gauss} with respect to the modulus $p$.
Cases I and III are handled in the same way.
